\chapter{Control Flow}
\index{control flow}
\label{ch:control-flow}

In most languages, control flow is managed with statements: \texttt{if}/\texttt{else} blocks, \texttt{for} loops, \texttt{switch} statements. In Scheme, all control flow is handled by expressions that produce values. This chapter covers \texttt{if}, \texttt{cond}, \texttt{case}, boolean combinators, \texttt{when}/\texttt{unless}, sequencing with \texttt{begin}, and iteration with \texttt{do}.

\section{\texttt{if}: The Fundamental Conditional}
\index{if}

The \texttt{if} expression is the most basic conditional. It has the form:

\begin{lstlisting}
(if test consequent alternate)
\end{lstlisting}

If \texttt{test} evaluates to any true value (anything that is not \texttt{\#f}), the \texttt{consequent} is evaluated and its value returned. Otherwise, \texttt{alternate} is evaluated and returned:

\begin{lstlisting}[style=repl]
> (if (> 3 2) "yes" "no")
"yes"
> (if (< 3 2) "yes" "no")
"no"
> (if (zero? 0) 'zero 'nonzero)
zero
\end{lstlisting}

Because \texttt{if} is an expression, you can use it anywhere a value is expected:

\begin{lstlisting}[style=repl]
> (+ 1 (if #t 10 20))
11
> (string-append "The answer is "
                 (if (even? 4) "even" "odd"))
"The answer is even"
\end{lstlisting}

\subsection{One-Armed \texttt{if}}

You can omit the alternate:

\begin{lstlisting}[style=repl]
> (if (> 5 3) (display "big\n"))
big
\end{lstlisting}

When the test is false and there is no alternate, the result is unspecified. In practice, use \texttt{when} (see below) for one-armed conditionals---it communicates intent more clearly.

\section{\texttt{cond}: Multiple Branches}
\index{cond}

\texttt{cond} is Scheme's equivalent of Python's \texttt{if}/\texttt{elif}/\texttt{else}. It tests multiple conditions in order:

\begin{lstlisting}[style=repl]
> (define (describe-temp celsius)
    (cond
      ((< celsius 0) "freezing")
      ((< celsius 10) "cold")
      ((< celsius 20) "mild")
      ((< celsius 30) "warm")
      (else "hot")))
> (describe-temp -5)
"freezing"
> (describe-temp 15)
"mild"
> (describe-temp 35)
"hot"
\end{lstlisting}

Each clause is \texttt{(test expr ...)}. The first clause whose test returns true has its expressions evaluated, and the last one's value is returned. The \texttt{else} clause matches if nothing else did.

\subsection{The \texttt{=>} Syntax}

A special clause form \texttt{(test => proc)} passes the test result to a procedure:

\begin{lstlisting}[style=repl]
> (cond
    ((assv 2 '((1 "one") (2 "two") (3 "three")))
     => cadr)
    (else "not found"))
"two"
\end{lstlisting}

Here, \texttt{assv} returns \texttt{(2 "two")} (a true value), and \texttt{=>} passes it to \texttt{cadr}, which extracts \texttt{"two"}. This avoids evaluating the lookup twice.

\section{\texttt{case}: Matching Against Values}
\index{case}

\texttt{case} compares a value against sets of constants, similar to a \texttt{switch} statement in C or \texttt{match} on literals in Python:

\begin{lstlisting}[style=repl]
> (define (day-type day)
    (case day
      ((monday tuesday wednesday thursday friday)
       "weekday")
      ((saturday sunday)
       "weekend")
      (else "unknown")))
> (day-type 'saturday)
"weekend"
> (day-type 'tuesday)
"weekday"
\end{lstlisting}

Each clause lists datum values in parentheses. The key is compared using \texttt{eqv?}\index{eqv?}, which tests equality for numbers, characters, booleans, and symbols (but not strings---two strings with the same content may not be \texttt{eqv?}). For matching against strings, use \texttt{cond} with \texttt{string=?} instead.

\begin{lstlisting}[style=repl]
> (case (+ 2 3)
    ((1 2 3) "small")
    ((4 5 6) "medium")
    (else "large"))
"medium"
\end{lstlisting}

\begin{note}[Coming from Python]
\texttt{case} is similar to Python 3.10's \texttt{match/case} but limited to literal value matching (no pattern destructuring). For pattern matching on structure, Scheme uses recursive functions or \texttt{cond} with predicates.
\end{note}

\section{Boolean Combinators: \texttt{and} and \texttt{or}}
\index{and}\index{or}

\texttt{and} and \texttt{or} are short-circuiting\index{short-circuit evaluation} expressions that return the determining value---not just \texttt{\#t} or \texttt{\#f}.

\subsection{\texttt{and}}

\texttt{and} evaluates its arguments left to right. If any is \texttt{\#f}, it returns \texttt{\#f} immediately. If all are true, it returns the last value:

\begin{lstlisting}[style=repl]
> (and 1 2 3)
3
> (and 1 #f 3)
#f
> (and (> 3 2) (< 5 10) "both true")
"both true"
> (and)
#t
\end{lstlisting}

\subsection{\texttt{or}}

\texttt{or} evaluates its arguments left to right. If any is true, it returns that value immediately. If all are \texttt{\#f}, it returns \texttt{\#f}:

\begin{lstlisting}[style=repl]
> (or #f #f 42)
42
> (or #f #f #f)
#f
> (or (memq 'b '(a b c)) 'not-found)
(b c)
> (or)
#f
\end{lstlisting}

\subsection{Using \texttt{or} for Defaults}

A common pattern is using \texttt{or} to provide fallback values:

\begin{lstlisting}[style=repl]
> (define (get-port config)
    (or (cond ((assoc "port" config) => cdr)
              (else #f))
        8080))
> (get-port '())
8080
> (get-port '(("port" . 3000)))
3000
\end{lstlisting}

If the \texttt{cond} finds a \texttt{"port"} entry, it extracts the value with \texttt{cdr}; otherwise it returns \texttt{\#f}. The \texttt{or} then provides the default when the result is \texttt{\#f}.

\subsection{Using \texttt{and} for Guarded Access}

\texttt{and} is useful for safe chaining when a value might be \texttt{\#f}:

\begin{lstlisting}[style=repl]
> (define data '(("name" . "Alice")))
> (and (assoc "name" data)
       (cdr (assoc "name" data)))
"Alice"
> (and (assoc "email" data)
       (cdr (assoc "email" data)))
#f
\end{lstlisting}

\begin{note}[Coming from Python]
Scheme's \texttt{and}/\texttt{or} behave identically to Python's. \texttt{(or val default)} is like Python's \texttt{val or default}. \texttt{(and x (f x))} is like \texttt{x and f(x)}.
\end{note}

\section{\texttt{not}}
\index{not}

\texttt{not} negates a boolean value. It returns \texttt{\#t} only when given \texttt{\#f}:

\begin{lstlisting}[style=repl]
> (not #f)
#t
> (not #t)
#f
> (not 0)
#f
> (not '())
#f
\end{lstlisting}

\begin{warning}[Zero is true in Scheme]
Since only \texttt{\#f} is false in Scheme, \texttt{(not 0)} is \texttt{\#f}---not \texttt{\#t}. This is the opposite of Python, where \texttt{not 0} is \texttt{True}. The same applies to the empty list, the empty string, and every other value: only \texttt{\#f} is false.
\end{warning}

\section{\texttt{when} and \texttt{unless}}
\index{when}\index{unless}

\texttt{when} is a one-armed conditional for side effects---it executes its body only if the test is true:

\begin{lstlisting}[style=repl]
> (define x 10)
> (when (> x 5)
    (display "x is large")
    (newline))
x is large
\end{lstlisting}

\texttt{unless} is the opposite---it executes when the test is false:

\begin{lstlisting}[style=repl]
> (unless (= x 0)
    (display "x is nonzero")
    (newline))
x is nonzero
\end{lstlisting}

Both forms accept multiple expressions in their body (implicit \texttt{begin}). Use them instead of one-armed \texttt{if} when you want to perform actions without needing a result value.

\section{\texttt{begin}: Sequencing}
\index{begin}

\texttt{begin} groups multiple expressions into one, evaluating them in order and returning the last value:

\begin{lstlisting}[style=repl]
> (begin
    (display "first\n")
    (display "second\n")
    42)
first
second
42
\end{lstlisting}

You rarely need explicit \texttt{begin} because most forms (\texttt{define} function bodies, \texttt{when}, \texttt{unless}, \texttt{cond} clauses) already accept multiple body expressions. It is mainly useful inside \texttt{if} when you need multiple expressions in one branch:

\begin{lstlisting}[style=repl]
> (if (> 5 3)
      (begin
        (display "yes\n")
        "positive")
      "negative")
yes
"positive"
\end{lstlisting}

\section{\texttt{do}: Iterative Loops}
\index{do}

Scheme has no dedicated \texttt{while} or \texttt{for} statement. Iteration is usually expressed with recursion, which Chapter~\ref{ch:functions} covers in depth. For loops with explicit loop variables, Scheme provides the \texttt{do} form:

\begin{lstlisting}
(do ((var init step) ...)
    (test result ...)
  body ...)
\end{lstlisting}

Each variable starts at \texttt{init} and is updated to \texttt{step} each iteration. The loop terminates when \texttt{test} becomes true, returning \texttt{result}:

\begin{lstlisting}[style=repl]
> (do ((i 0 (+ i 1))
       (sum 0 (+ sum i)))
      ((= i 10) sum))
45
\end{lstlisting}

Here is \texttt{do} building a reversed list:

\begin{lstlisting}[style=repl]
> (do ((lst '(1 2 3 4 5) (cdr lst))
       (result '() (cons (car lst) result)))
      ((null? lst) result))
(5 4 3 2 1)
\end{lstlisting}

\begin{note}[Coming from Python]
\texttt{do} is roughly equivalent to a Python \texttt{for} or \texttt{while} loop. Most Scheme programmers prefer recursion (Chapter~\ref{ch:functions}) or named \texttt{let} (Chapter~\ref{ch:bindings}) for loops---but \texttt{do} needs nothing beyond this chapter, and it shines when several loop variables change together.
\end{note}

\section{Putting It All Together}

Here is a function that demonstrates several control flow forms working together:

\begin{lstlisting}
(define (fizzbuzz n)
  (do ((i 1 (+ i 1)))
      ((> i n))
    (cond
      ((and (zero? (modulo i 3))
            (zero? (modulo i 5)))
       (display "FizzBuzz"))
      ((zero? (modulo i 3))
       (display "Fizz"))
      ((zero? (modulo i 5))
       (display "Buzz"))
      (else
       (display i)))
    (display " ")))
\end{lstlisting}

\begin{lstlisting}[style=repl]
> (fizzbuzz 15)
1 2 Fizz 4 Buzz Fizz 7 8 Fizz Buzz 11 Fizz 13 14 FizzBuzz
\end{lstlisting}

This uses \texttt{do} for the loop, \texttt{cond} for the multi-branch logic, and \texttt{and} to combine conditions.

\section*{Summary}

\begin{itemize}
    \item \texttt{if} is the fundamental two-way conditional; \texttt{cond} and \texttt{case} handle multiple branches.
    \item \texttt{and} and \texttt{or} short-circuit and return the determining value, which is handy for defaults and guards.
    \item \texttt{when} and \texttt{unless} run a body under a single condition with no else branch.
    \item \texttt{begin} sequences expressions, and \texttt{do} expresses iterative loops.
    \item Because conditionals are expressions, they produce values rather than merely directing control.
\end{itemize}

\section*{Exercises}

\begin{exercise}[Configurable FizzBuzz]
The chapter showed a standard FizzBuzz implementation. Generalize it: write a function \code{fizzbuzz-custom} that takes an upper bound \code{n} and an association list of \code{(divisor . word)} pairs. For each number from 1 to \code{n}, print the concatenation of all matching words, or the number itself if no divisor matches:

\begin{lstlisting}[style=repl]
> (fizzbuzz-custom 15 '((3 . "Fizz") (5 . "Buzz")))
1 2 Fizz 4 Buzz Fizz 7 8 Fizz Buzz 11 Fizz 13 14 FizzBuzz
> (fizzbuzz-custom 10 '((2 . "Even") (3 . "Tri")))
1 Even Tri Even 5 EvenTri 7 Even Tri Even
\end{lstlisting}

Use \code{do} for the outer loop and a helper function to build the output for one number. This exercises \code{cond}, \code{when}, \code{do}, and string operations together.
\end{exercise}

\begin{exercise}[Letter Grades]
Write a function \code{grade} that takes a numeric score (0--100) and returns a letter grade as a string: \code{"A"} for 90+, \code{"B"} for 80--89, \code{"C"} for 70--79, \code{"D"} for 60--69, and \code{"F"} below 60. Use \code{cond}:

\begin{lstlisting}[style=repl]
> (grade 95)
"A"
> (grade 73)
"C"
> (grade 42)
"F"
\end{lstlisting}

Bonus: add input validation---return \code{"Invalid"} if the score is not between 0 and 100.
\end{exercise}

\begin{exercise}[Clamping a Value]
Write a function \code{clamp} that takes three arguments---\code{value}, \code{lo}, and \code{hi}---and returns \code{value} restricted to the range [\code{lo}, \code{hi}]. Use \code{cond}:

\begin{lstlisting}[style=repl]
> (clamp 5 1 10)
5
> (clamp -3 0 100)
0
> (clamp 150 0 100)
100
\end{lstlisting}

In Python, this is \texttt{max(lo, min(hi, value))}. How does the \code{cond} version compare for readability?
\end{exercise}

In the next chapter, we turn to functions in depth---including recursion, the technique Scheme programmers reach for where other languages use loops.

\chapterend
